%% Chapter 4 - Product and Solution Design
%% Planned sections (master structure): 4.1 Requirements; 4.2 Component
%% architecture within the complete solution; 4.3 Design alternatives and
%% decisions; 4.4 Detailed design; 4.5 Standards and professional practice.
\chapter{Product and Solution Design}\label{ch:product-design}

\section{Requirements}\label{sec:requirements}

\subsection*{Product context}

NeuroShell IRE (the Intent Recognition Engine) is Component~1 of the
NeuroShell Core Framework, a language-driven security-automation stack. The
product is a domain-specific natural-language understanding gateway: it
converts a free-text offensive-security command into a validated,
schema-constrained intent contract that downstream components can consume
without ever seeing free-form model output (Section~\ref{sec:requirements}).
The contract-first boundary is the central requirement of the product: every
output exposed to the toolchain must be machine-validated, and no unvalidated
LLM text may reach it (Chapters~\ref{ch:introduction} and~\ref{ch:methodology}).

The IRE runs as a local daemon on the operator's device, serving other
components and human operators over HTTP. The complete NeuroShell framework
places two further safety gates downstream of the IRE—a planner that enforces
its own RFC~1918 and dangerous-pattern checks, and an executor that runs
recovery-driven command execution—so the IRE is designed as the first, but not
the only, enforcement point.

\subsection*{User classes}

The product serves four human role classes in a strict privilege hierarchy
(viewer, analyst, operator and admin) and two classes of machine consumer
(downstream services and human-in-the-loop reviewers). Table~\ref{tab:req-roles}
summarises the roles; the full intent-permission matrix and behavioural rules
are given in Section~\ref{sec:detailed-design}.

\begin{table}[h]
\centering
\caption{Role model for the IRE. The default role is analyst; unknown roles
coerce to viewer (least privilege).}\label{tab:req-roles}
\begin{tabular}{@{}lll@{}}
\toprule
\textbf{Role} & \textbf{Hierarchy} & \textbf{Privilege character}\\
\midrule
viewer & 0 (least) & Observation only; pre-inference fast-fail for any intent\\
analyst & 1 (default) & Passive recon, network scan and vulnerability audit\\
operator & 2 & All intents, including the three sensitive classes\\
admin & 3 (most) & All intents plus model, policy and feature management\\
\bottomrule
\end{tabular}
\end{table}

\subsection*{Functional requirements}

The functional requirements are grouped under the product functions verified
in the implementation (Section~\ref{sec:component-architecture}):
\begin{itemize}
\item \textbf{Core understanding.} Normalise the input utterance; resolve
offensive-security aliases (for example \texttt{EternalBlue} to
CVE-2017-0144) to canonical identifiers; classify the request into one of nine
top-level intents; and refine it into a hierarchical sub-intent that drives
downstream tool selection.
\item \textbf{Safety and security.} Detect adversarial input before inference;
enforce role-based access control both before and after inference; validate
the model's output against typed schema, regular-expression, network and
engagement-scope rules; classify every failed validation into one of six
hallucination classes; and emit an audit record for every parse.
\item \textbf{Intelligence and quality.} Maintain per-session context across
turns; serve identical intents from a semantic cache; gate low-confidence
decisions through a self-consistency stage; and expose an optional
XAI attribution path.
\item \textbf{System and integration.} Expose a REST API for parsing, execution
and administration; push validated contracts to downstream components on a
feature-flagged basis; persist an append-only audit trail; and make every
guard individually disable-able through configuration.
\end{itemize}

\subsection*{Non-functional requirements}

\begin{itemize}
\item \textbf{Performance.} Adversarial detection (M1) must complete in low
milliseconds; the pre-inference RBAC fast-fail runs before any model cost; and
cache hits must short-circuit the pipeline (Section~\ref{sec:measures}).
\item \textbf{Reliability.} A failure in non-critical middleware (session
context, cache, adversarial detection) must degrade gracefully, not abort the
request.
\item \textbf{Security.} The service holds no credentials beyond an operator
API key; rate limiting protects the model backend; and no free-form model
output is ever returned to a caller (Section~\ref{sec:scope-limitations-assumptions}).
\item \textbf{Operability.} Every stage and middleware is individually
configurable via feature flags; the model family and engagement scope are
administratively reconfigurable at runtime; and the whole evaluation is
reproducible from a single manifest (Section~\ref{sec:reproducibility}).
\end{itemize}


\section{Component Architecture within the Complete Solution}\label{sec:component-architecture}

\subsection*{Position in the NeuroShell framework}

The IRE is Component~1 of a four-component stack. Component~2 (the Dynamic
Planner) consumes \texttt{PlannerContract} objects to synthesise an execution
plan; Component~3 (AEERE) executes the plan with adaptive error recovery; and
Component~4 (AVAE) performs post-execution vulnerability analysis. Downstream
integration is optional and feature-flagged
(\texttt{planner\_enabled}, \texttt{component3\_enabled},
\texttt{component4\_enabled}), so the IRE is a complete product on its own and
a drop-in first stage of the framework. Machine consumers are authenticated by
a shared API key and receive only validated schemas; they never receive
free-form LLM output.

\subsection*{Top-level architecture}

Table~\ref{tab:stage-modules} maps the verified pipeline stages to their
implementing modules under \texttt{src/}. A deterministic core of seven
serialised stages (S1--S7, where S6B is the network sub-validator of S6 rather
than an eighth stage) transforms the raw utterance into a validated intent,
while five middleware modules enforce mandatory safety and quality properties
around the core. The guarded decision function and safety invariant
(Equations~\ref{eq:guarded-decision} and~\ref{eq:safety-invariant},
Chapter~\ref{ch:methodology}) are realised by exactly these components.

\begin{table}[h]
\centering
\small
\caption{Verified pipeline stages and middleware under \texttt{src/}.}\label{tab:stage-modules}
\begin{tabular}{@{}p{2.9cm}>{\raggedright\arraybackslash}p{4.7cm}>{\raggedright\arraybackslash}p{5.8cm}@{}}
\toprule
\textbf{Component} & \textbf{Module} & \textbf{Primary responsibility}\\
\midrule
S1 InputNormalizer & \texttt{preprocessing/}\allowbreak\texttt{input\_}\allowbreak\texttt{normalizer.py} & NFKC normalisation, case folding, jargon standardisation\\
S2 AliasResolver & \texttt{preprocessing/}\allowbreak\texttt{alias\_}\allowbreak\texttt{resolver.py} & Alias-to-canonical (CVE, tool, action) resolution\\
S3 OllamaInference & \texttt{inference/}\allowbreak\texttt{ollama\_}\allowbreak\texttt{inference\_}\allowbreak\texttt{engine.py} & Local LLM inference with strict-JSON directive\\
S4 JSONParser & \texttt{validation/}\allowbreak\texttt{json\_}\allowbreak\texttt{parser.py} & JSON extraction; strips fences and thinking tags\\
S5 SchemaValidator & \texttt{validation/}\allowbreak\texttt{schema\_}\allowbreak\texttt{validator.py} & Typed \texttt{IntentSchema} enforcement\\
S6 RegexValidator & \texttt{validation/}\allowbreak\texttt{regex\_}\allowbreak\texttt{validator.py} & IP / CIDR / CVE / port pattern verification\\
S6B NetworkValidator & \texttt{validation/}\allowbreak\texttt{network\_}\allowbreak\texttt{validator.py} & Scope containment, CIDR sanity, target-type consistency\\
S7 ScopeGuard & \texttt{validation/}\allowbreak\texttt{scope\_}\allowbreak\texttt{guard.py} & RFC~1918 enforcement and injection re-scan\\
M1 AdversarialDetector & \texttt{middleware/}\allowbreak\texttt{adversarial\_}\allowbreak\texttt{detector.py} & Three-layer threat scoring, block/flag before inference\\
M2/M5 RbacGuard & \texttt{middleware/}\allowbreak\texttt{rbac\_}\allowbreak\texttt{guard.py} & Pre-inference fast-fail and post-inference intent check\\
M3 SessionContext & \texttt{middleware/}\allowbreak\texttt{session\_}\allowbreak\texttt{context.py} & Multi-turn context, last-target carry-over, TTL expiry\\
M4 SemanticCache & \texttt{middleware/}\allowbreak\texttt{semantic\_}\allowbreak\texttt{cache.py} & Exact-hash and embedding-similarity cache hit\\
Sub-intent classifier & \texttt{middleware/}\allowbreak\texttt{sub\_}\allowbreak\texttt{intent\_}\allowbreak\texttt{classifier.py} & Refines intent into hierarchical sub-intents\\
ContractBuilder & \texttt{pipeline/}\allowbreak\texttt{contract\_}\allowbreak\texttt{builder.py} & Synthesises the downstream \texttt{PlannerContract}\\
\bottomrule
\end{tabular}
\end{table}

\subsection*{Processing order}

Middleware runs around the core: M1 classifies the raw input, M2 fast-fails
restricted roles before any model cost, M3 injects session context, M4 serves
cache hits that short-circuit the LLM entirely, and the core (S1--S7, S6B) then
runs in strict sequence. Following the core, M5 re-validates the detected
intent against the role matrix, and the ContractBuilder produces the
downstream contract. The appendix and the threat model
(Sections~\ref{sec:gap-synthesis} and~\ref{sec:research-design}) treat the
downstream boundary: the IRE renders no free-form content and retrieves
nothing, closing the indirect-injection and tool-injection channels.


\section{Design Alternatives and Decisions}\label{sec:design-alternatives}

\subsection*{Local, offline decoding versus a hosted model}

The first architectural decision is the decoder. A cloud-hosted model offers
higher raw accuracy and easier scaling, but contradicts the product
requirement that no request leaves the device and no credentials be held
(Sections~\ref{sec:requirements} and~\ref{sec:ethics}). The chosen design runs
a small model locally through Ollama in CPU-only mode, which reflects the
intended field deployment and keeps the trust boundary on a single device. Two
backends are kept configurable so generality can be measured
(Chapter~\ref{ch:methodology}).

\subsection*{Structured contract versus free-form prose}

Raw LLM output is plausible prose and a machine-dispatched action must not be
built from prose. The chosen design forces the model to emit a fixed JSON
schema (S3) and then validates it through S4--S7 before anything is returned.
Only an output that survives these stages is exposed; otherwise a
stage-specific classified error is returned. This converts a plausibility claim
into a verifiable one and is the decision behind
Equations~\ref{eq:guarded-decision} and~\ref{eq:safety-invariant}.

\subsection*{Deterministic rules versus model self-checking}

Two validation philosophies are available: ask the model to review its own
output, or recompute the checks with deterministic rules. Self-checking is
drift-prone and reuses the same faulty estimator. The chosen design recomputes
every syntactic, network and scope check with deterministic code (S5--S7,
S6B) that depends on no further LLM calls. This is why the deterministic guard
path is bit-exact and auditable (Section~\ref{sec:reproducibility}).

\subsection*{No retrieval versus retrieval-augmented generation}

Retrieval-augmented generation is a popular way to ground answers, but
retrieved text is itself a prompt-injection vector~\cite{greshake2023indirect},
and a retrieval channel expands the offline trust boundary. The chosen design
performs no retrieval at inference time; offline CVE grounding is loaded at
startup only. This decision is what makes the threat model closed
(Section~\ref{sec:deployment-properties}).

\subsection*{Two-phase RBAC versus a single check}

Role enforcement can be a single post-hoc check or a repeated invariant. The
chosen design enforces RBAC twice: a pre-inference fast-fail (M2) that rejects
restricted roles before any model cost, and a post-inference check (M5) that
validates the final intent against the role matrix. A single bypassed guard is
therefore not sufficient to manufacture an authorised action.

\subsection*{Self-consistency versus single-pass confidence}

The model supplies a confidence score, but calibrated scores are unreliable for
small offline models. The chosen design gates low-confidence decisions through
a self-consistency stage that samples further responses and measures agreement,
accepting the bounded latency cost only when triggered. This is the
sampling-based disagreement intuition applied to a single structured
classification rather than to free-form text~\cite{wang2023selfconsistency}.

\subsection*{Semantic caching}

Identical intents recur in operator workflows. The chosen design adds a
two-channel semantic cache (exact hash and embedding similarity) that
short-circuits the LLM on hits, cut latency from seconds to tens of
milliseconds per hit, and improves determinism; the cache is invalidated when
the offline CVE-grounding database is reindexed. The latency figures are not
quoted here pending author confirmation (Section~\ref{sec:measures}).


\section{Detailed Design}\label{sec:detailed-design}

\subsection*{Core stages (S1--S7, S6B)}

S1 normalises casing, whitespace and Unicode (NFKC) and standardises jargon
terms into canonical phrases. S2 resolves aliases such as \texttt{nmap -sS}
to a SCAN action and exploit surnames to canonical CVE identifiers via
\texttt{data/alias\_map.json}. S3 instructs the model to emit a fixed JSON
schema, at the sampling configuration pinned in the evaluation manifest
(Section~\ref{sec:reproducibility}). S4 recovers the JSON payload, stripping
code fences, markdown artefacts and thinking tags. S5 enforces the typed
\texttt{IntentSchema} contract: exactly one of nine \texttt{IntentType} values,
a matched target type, ports in 1--65535, CVE syntax \texttt{CVE-YYYY-NNNN},
confidence in \lbrack 0,1\rbrack, and the invariant that \texttt{AMBIGUOUS}
requires confidence below 0.5 while \texttt{REJECTED} requires a
\texttt{rejection\_reason}. S6 deterministically re-verifies every entity the
model produced (IPv4/IPv6, CIDR, CVE, port, URL) and the absence of shell
metacharacters. S6B applies engagement-scope and topology rules: targets must
fall inside the configured engagement CIDR (default 192.168.0.0/16), the
widest permitted scan prefix is 30 and the narrowest 32, and type/value
mismatches fail here, generating \texttt{ARCH\_WARNING} and
\texttt{SCOPE\_WARNING} entries. S7 is the final enforcement gate: it rejects
out-of-scope targets in strict mode, warns in research mode, applies RFC~1918
membership~\cite{rfc1918}, and re-scans the raw input as an injection gate.

\subsection*{Middleware}

M1 is a three-layer adversarial detector operating before inference: a lexical
layer scores known injection phrases and control tokens (weights 0.8--1.0), a
structural layer scores encoded payloads, script-like constructs and
imperative-but-unclassifiable phrasing (0.3--0.9), and a scoring layer fuses
the signals; a cumulative score at or above 0.95 auto-blocks immediately, and
0.7 flags the request for downstream treatment (Section~\ref{sec:procedure}).
M2/M5 implement RBAC in two phases around the core as described in
Section~\ref{sec:design-alternatives}; the policy lives in
\texttt{config/rbac\_policy.py} and rules that unknown roles coerce to
\texttt{viewer}, the least-privileged class. M3 maintains per-session turn
history with a TTL, injects the recent memory window into the prompt and
supports anaphoric follow-ups such as ``the same target''. M4 serves exact-hash
and cosine-similarity cache hits, evicts LRU-style, and is invalidated with the
CVE-grounding hash. The sub-intent classifier refines the top-level intent into
a hierarchical sub-intent (for example \texttt{NETWORK\_SCAN.SYN\_STEALTH}),
which the ContractBuilder maps to a suggested tool and flags
(Section~\ref{sec:contract-generation}).

\subsection*{Hallucination taxonomy and enforcement}

Every failed validation is classified into one of six hallucination classes via
the enforcement-policy configuration module
(\texttt{config/enforcement\_policy.py}) and becomes an audit-grade
\texttt{ValidationFinding}. The policy maps classes to
BLOCK/WARN/ESCALATE outcomes: fabricated parameters, contradictory
action-target pairings, target-type mismatches and malformed fabricated CVEs
hard-block; ungrounded confidence (syntactically valid but unindexed CVEs)
warns; and out-of-scope targets warn in research mode but block in strict
production mode. Escalations route through the admin escalation queue for
explicit human approval.

\subsection*{Contract generation and administration}\label{sec:contract-generation}

The ContractBuilder synthesises a downstream-ready \texttt{PlannerContract}
with a strict, fully validated field set (intent, sub-intent, validated target,
ports, modifiers, CVE IDs, tool hint, schedule, confidence, rejection reason
and session ID), a human-readable description, a priority level (scheduled,
critical for exploitation, high for CVE audits, medium otherwise) and a
generated suggested command. Exploitation and password-attack flows are marked
as requiring explicit human approval. Management is exposed through the REST
API: parsing, execution, metrics, sessions, and admin views for RBAC, feature
flags, models, enforcement policy, escalations, audit summary and the
engagement scope. An append-only audit logger persists a SHA-256 hashed form
of operator content with the intent, role, latency and outcome as structured
fields.

\subsection*{Deployment properties}\label{sec:deployment-properties}

The service runs headless, CPU-only, retrieving nothing at inference time and
holding no credentials beyond the operator API key; its only persistent assets
are the model weights, the authorization policy, the session context and the
audit database. Because nothing is retrieved and no free-form output is
rendered downstream, the two principal remote attack classes (indirect prompt
injection via retrieved data~\cite{greshake2023indirect} and tool injection via
unfiltered text) are structurally reduced to classifying a single user
utterance.

\subsection*{Boundaries of the design record}

The audit rule for this chapter is that components are described only when
their implementing module is present under \texttt{src/}. Four capabilities that
appear in the project documentation are therefore flagged rather than
described as implemented: \texttt{spaCy} text pipelining is referenced in the
SRS dependencies but has no module in \texttt{src/}; the \texttt{jsonschema}
runtime is a documented dependency while the actual validator in
\texttt{src/} uses Pydantic models; LoRA fine-tuning exists as a configuration
file (\texttt{config/lora\_config.yaml}) but is planned Phase~2 work, not an
implemented execution path; Mahalanobis anomaly detection and XAI token
attribution appear as configuration (feature-flag \texttt{xai: false}) and
tuning parameters, not as wired components. Where these names recur in
discussion, they are treated as unimplemented design candidates.


\section{Standards and Professional Practice}\label{sec:standards}

The design follows established standards where they apply. The SRS was written
against the IEEE~830 recommended practice for software requirements
specifications, and the architecture's private-address enforcement is grounded
in RFC~1918~\cite{rfc1918}. The threat catalogue that motivates the guard
stack follows the OWASP Top~10 for LLM applications~\cite{owasp2025}, and
offline CVE grounding draws on the National Vulnerability Database as a
canonical identifier source~\cite{nistnvd}. Security tooling references, such
as the suggested scanning tools and the exploitation framework, follow the
documented syntax of the respective vendors~\cite{lyon2009nmap,metasploit} so
that generated contracts are executable by the downstream components as
intended.

Professional practice in this domain is primarily the ethics of dual use
(Section~\ref{sec:ethics}): the system can drive offensive tooling, so the
design gates every such action behind role authorisation, declared engagement
scope and explicit human approval for the most sensitive flows, and restricts
targets to operator-declared private ranges~\cite{rfc1918}. Responsible-AI
practice is observed by treating the model as an untrusted proposition
generator whose output is validated deterministically and reported honestly
(Sections~\ref{sec:data-analysis} and~\ref{sec:ethics}); capability, limitation
and residual failure modes are published rather than suppressed. Accessibility
is respected at the interface level (clear, structured API responses rather
than prose), and because the product processes no personal data, it falls
outside personal-data protection regimes; the applicable disclosures and
declarations are collated in Chapter~\ref{sec:ethics} and
Appendix~\ref{app:b}.